% ------------------------------------------------------------
% Page 93
% ------------------------------------------------------------

\sectiondate{Les fêtes des Oubrets}{}

Il ne m’est pas possible de passer sous silence la célébration périodique d’une vraie
fête trimestrielle de cette paroisse. Je n’ai rien connu de semblable ailleurs ! Chaque
fois que le mois comportait cinq dimanches, c’est à dire à quatre reprises dans
l’année, les cultes habituels étaient supprimés à Meyrueis et à Gatuzieres. On affichait
« relâche ». Le pasteur, accompagné parfois de deux ou trois personnes, montait
jusqu’aux Oubrets. Ce hameau, à 1.000m d’altitude, terminait la route partie de
Meyrueis en un cul de sac, proche de la source de la Brèze. Il comportait alors un
certain nombre de maisons à l’aspect délabré et groupait une population relativement
jeune et dense.

On y remarquait le petit temple protestant que l’Association Cultuelle prêtait à la
commune de Meyrueis afin de servir d’école. Pour en conserver la propriété, d’après
la convention de la loi de séparation des Eglises et de l’Etat, on devait y organiser une
manifestation quelconque (culte, congrès, concert...) au moins une fois par an. D’où
la décision de ce culte trimestriel. Les enfants de ce lieu, tous protestants,
fréquentaient la classe dirigée par une institutrice dévouée, habitant un peu plus bas à
Pourcarès. Le cinquième dimanche du mois, donc, s’y déroulait une fête au village.

On la préparait avec beaucoup de soins. En premier lieu le culte. Personne n’aurait
songé à le manquer. Puis un grand repas en commun, présidé par le pasteur et sa
femme. Enfin l’après-midi et même le début de la soirée, toujours dans le même
local, se prolongeait la rencontre par une sorte de veillée joyeuse et fraternelle. On
mangeait ensemble les produits de la pêche et de la chasse, truites et gibiers, ainsi que
de l’élevage des fermes de l’endroit. Les coups de fusil claquaient dès l’aube, en
saison et hors de saison (parfois en contrebande). Un lapin, un lièvre, des perdreaux
en faisaient les frais.

L’auto paroissiale pétaradante troublait le silence de la nature. Elle chauffait sur ce
chemin, à peine carrossable, non goudronné, aux montées excessives. J’ai toujours eu
des difficultés pour atteindre le sommet, l’engin n’était pas approprié à ce parcours.
En hiver, le froid devenait intense, le village des Oubrets se trouvant en flanc de
montagne à proximité de la crête de l’Aigoual, mais sur la face Nord.

Quand je jugeai que le radiateur, bouillant comme une marmite surchauffée n’avait
plus d’eau, je m’arrêtais. J’allais chercher de gros glaçons dans le ruisseau tout
proche. Mon beau-frère Elie Brée, tout jeune alors, montant avec nous un jour, fit le
va et vient avec de la glace plein les mains : il participa à l’opération. Il m’en parle
encore ! On arrivait ainsi à rafraîchir un peu la voiture, au risque de faire éclater les
culasses. Mais une fois le bout de nos peines atteint, la journée se passait au mieux,
dans l’atmosphère la plus champêtre et la plus détendue.

% ------------------------------------------------------------
% Page 94
% ------------------------------------------------------------

Malgré les grands feux de bois on ne pouvait faire remonter suffisamment le
thermomètre. Les mamans n’avaient d’autre alternative, pour réchauffer leurs enfants,
que de les mettre au lit, en plein jour, sous les édredons, quand on en possédait. Un
fois remis en état ces enfants pouvaient retourner dehors à leurs jeux, quitte à revenir
sous les couvertures lorsque le gel les atteignait à nouveau.

Véritablement se vivait là, une vie primitive, presque animale, faite de dépouillement,
de souffrance et de pauvreté. C’était semblable à ce que nous avions vu au Pompidou
dans certaines fermes.

Les pouvoirs publics consentaient un effort financier en faveur de ce hameau. Chaque
famille touchait annuellement une petite indemnité de maintien. Quand l’Inspecteur
des Eaux et Forêts passait pour compter le nombre exact de foyers intéressés par cette
mesure, on allumait des feux dans toutes les maisons abandonnées. Il s’agissait de
bénéficier au maximum de cette aide de l’Etat pour lutter contre la dépopulation des
montagnes. Le total à partager restait cependant bien modeste, et l’encouragement à
occuper les lieux bien relatif. On l’appréciait quand même !

Le culte et la fête apportaient à ces « braves gens » comme un rayon de soleil et
transformaient momentanément, en tous cas, leur horizon. Nous étions tous réjouis et
réconfortés de ce contact avec ces humbles paysans déshérités, mais qui ne se
plaignaient pas de leur sort. Au contraire, la satisfaction éclatait sur leurs visages, et
je ne crois pas que, en temps ordinaire, ils aient connu vraiment la tristesse ou
l’amertume.

Tous les plats qui constituaient le repas provenaient du dur labeur (y compris la pêche
ou la chasse) de ces hommes et de ces femmes sauf un : le plat de morue servi dans la
sauce blanche légèrement colorée. Je n’ai jamais rien dégusté de plus savoureux et
aujourd’hui je donnerai beaucoup pour en retrouver le goût. Voilà une recette
culinaire précieuse qui s’est probablement perdue... ! (mais non, vous trouverez celle de
Marianne Loupiac à la fin de ces pages)

Et puis cette vallée de la Brèze, relativement étroite en ses débuts, avec ses grands prés
séparant la route de la rivière, ses arbres majestueux, hêtres et érables, ses tournants
continuels suivis parfois d’une longue ligne droite, vallée de plus en plus abrupte à
mesure qu’on la pénètre, jusqu’à son dégagement par des lacets au bout desquels se
dévoile brusquement l’agglomération des Oubrets : quelle merveille de la nature ! Je
l’ai parcourue à nouveau récemment, confondu par ce spectacle.

\hfill Jean Gounelle
